\chapter{Testing and Evaluation}
\label{chap:testing}

This chapter presents the quantitative results of running the pipelines described in
Chapter~\ref{chap:implementation}.
The goal is to determine whether injecting parameters from the Asset Administration Shell causes
the simulated robot to behave more like the physical UR3.
Results are organised by metric: execution timing, joint position accuracy, TCP path deviation,
and motor current.
A summary table at the end consolidates all findings and identifies what improved, what did not,
and why.


\section{Evaluation Approach}
\label{sec:eval_methodology}

The study follows a quasi-experimental design \cite{cook1979quasi, wieringa2014design} in
which the physical robot serves as the ground-truth reference and the simulation variant
acts as the treatment condition (see Section~\ref{sec:research_methodology} for the full
methodological rationale).

Three pipelines are defined in this project.
Pipeline~1 (\textit{real}) provides the physical reference signal.
Pipeline~2 (\textit{sim\_aas}) is the primary AAS-enabled digital twin: URSim with payload
mass and TCP offset injected live from the BaSyx AAS server before each run.
Pipeline~3 is the AAS Python Simulator, a custom Python-based simulation model designed to
consume the full AAS submodel including the \textit{DynamicsParameters} submodel
(friction coefficients, damping factors, control latency, and current noise level).
Due to time constraints, Pipeline~3 was not completed within the project period; its design
is presented in Section~\ref{sec:aas_python_sim} and it is identified as the principal
direction for future work.

This chapter therefore evaluates one primary pairwise comparison: Pipeline~1 versus
Pipeline~2.
The \texttt{comparator.py} module aligns the two runs by interpolating both onto a shared
time axis, then computes three metrics:

\begin{itemize}
    \item \textbf{Per-joint RMSE} -- root-mean-square difference in joint angle (radians and
          degrees) between the simulation and the real robot at every time step.
    \item \textbf{TCP RMS deviation} -- RMS of the 3D Euclidean distance between the two runs'
          tool-centre-point positions (millimetres).
    \item \textbf{RMS joint current} -- root-mean-square motor current per joint (amperes).
\end{itemize}

These three metrics together cover kinematic fidelity (joint RMSE, TCP deviation) and
dynamic fidelity (current), giving a complete picture of what the AAS parameter injection
does and does not change.


\section{Dataset}
\label{sec:eval_dataset}

All runs were recorded to the shared Supabase database. Table~\ref{tab:dataset_overview}
summarises what was collected across the full project.

\begin{table}[htbp]
\centering
\begin{tabular}{lrr}
\toprule
Pipeline                       & Total runs recorded & Evaluation runs \\
\midrule
\textit{real} (Pipeline~1)     & 26                  & 6 \\
\textit{sim\_aas} (Pipeline~2) & 14                  & 1 \\
\bottomrule
\end{tabular}
\caption{Runs recorded for the evaluated pipelines. The evaluation subset is the set of
         full-length runs used for quantitative comparison.}
\label{tab:dataset_overview}
\end{table}

The six real evaluation runs completed in 110.01--110.03\,s (mean 110.02\,s,
$\sigma = 0.01$\,s), confirming stable and repeatable trajectory execution on the physical
robot.
This consistency justifies using a single representative real run as the reference signal.

The \textit{sim\_aas} evaluation run used the following parameters, fetched live from the BaSyx
server immediately before execution: payload mass $m = 1.2$\,kg at the tool origin, and TCP
offset $[-1.2\,\text{mm},\;0.64\,\text{mm},\;241.1\,\text{mm}]$ with a $90^{\circ}$
rotation about the X-axis.


\section{Execution Timing}
\label{sec:timing_analysis}

Before comparing accuracy, it is important to understand how the two pipelines differ in
execution speed, because a timing mismatch between runs directly affects how RMSE is computed.

\begin{figure}[H]
    \centering
    \includegraphics[width=0.78\textwidth]{Projekt/Figures/7_Test and Evaluation/fig8_cycle_time_bar.pdf}
    \caption{Measured cycle time (home $\to$ pick $\to$ place $\to$ home) per pipeline.
             Error bars show the standard deviation across evaluation runs.}
    \label{fig:cycle_time}
\end{figure}

\begin{table}[htbp]
\centering
\caption{Timing summary for the two representative evaluation runs.}
\label{tab:timing}
\begin{tabular}{lccc}
\toprule
Pipeline                       & Total duration (s) & Mean cycle time (s)  & Samples \\
\midrule
\textit{real} (Pipeline~1)     & 110.01             & $34.67 \pm 0.02$     & 13{,}752 \\
\textit{sim\_aas} (Pipeline~2) & 139.07             & $39.84 \pm 0.04$     & 17{,}385 \\
\bottomrule
\end{tabular}
\end{table}

Figure~\ref{fig:cycle_time} and Table~\ref{tab:timing} show two distinct timing effects.

Pipeline~2 is approximately 5\,s slower per cycle than the real robot; each cycle takes
approximately 39.8\,s in URSim versus 34.7\,s on the physical robot.
The physical UR3 runs on a dedicated real-time kernel with hardware-level motion interpolation.
URSim processes RTDE commands in a software loop, which adds latency to every motion
instruction.
This overhead is inherent to URSim and is independent of the AAS layer.

Pipeline~2 also carries an additional 14.4\,s dwell before the first trajectory cycle.
Before the first trajectory command, the pipeline must connect to the BaSyx server, retrieve
the submodel data, and apply the parameters via RTDE (see Figure~\ref{fig:pipeline2_sequence}
in Chapter~\ref{chap:implementation}).
Additionally, the \textit{sim\_aas} run begins recording from the URSim boot pose rather than
from the canonical home position, so the robot also executes a short move-to-home motion at
the start of the recording.
The combined result is that Pipeline~2 does not depart from the home position toward the first
pick waypoint until $t \approx 14.4$\,s, compared with $t \approx 1.85$\,s for the real run.

To ensure phase-coherent comparison, the comparator is configured to align each run at its
\emph{trajectory onset}: the first sample at which the robot departs the canonical home
position toward waypoint~1.
Pre-onset data (boot-to-home motion and AAS initialisation dwell) are excluded before computing
RMSE and TCP deviation.
Both runs are therefore compared over the same 108\,s trajectory window, starting from the
moment the arm begins moving toward the pick position.


\section{Joint Position Accuracy}
\label{sec:joint_accuracy}

\begin{figure}[H]
    \centering
    \includegraphics[width=0.88\textwidth]{Projekt/Figures/7_Test and Evaluation/fig8_joint_rmse_bar.pdf}
    \caption{Per-joint position RMSE between Pipeline~2 (\textit{sim\_aas}) and the real
             robot (trajectory-onset aligned).}
    \label{fig:joint_rmse}
\end{figure}

Table~\ref{tab:joint_rmse} gives the exact RMSE values for all six joints, and
Figure~\ref{fig:joint_rmse} shows them as a bar chart.

\begin{table}[htbp]
\centering
\begin{tabular}{lrr}
\toprule
Joint             & RMSE (rad) & RMSE ($^\circ$) \\
\midrule
J1 -- Base        & 0.1522 &  8.7 \\
J2 -- Shoulder    & 0.1754 & 10.1 \\
J3 -- Elbow       & 0.6607 & 37.9 \\
J4 -- Wrist~1     & 0.5318 & 30.5 \\
J5 -- Wrist~2     & 0.0325 &  1.9 \\
J6 -- Wrist~3     & 0.2403 & 13.8 \\
\midrule
Combined RMS      & 0.3723 & 21.3 \\
\bottomrule
\end{tabular}
\caption{Per-joint position RMSE between Pipeline~1 (\textit{real}) and Pipeline~2
         (\textit{sim\_aas}), trajectory-onset aligned.}
\label{tab:joint_rmse}
\end{table}

The largest errors appear at J3 and J4, which together account for roughly 75\,\% of the
total RMSE.
These are the joints with the largest angular excursions during the pick-and-place motion,
so systematic differences between the real robot's acceleration profile and URSim's
software-loop motion execution accumulate to significant angular offsets.

These errors reflect the inherent gap between URSim's software-loop motion execution and the
physical controller's real-time interpolation, not a deficiency of the AAS approach.
The robot executes trajectories via \texttt{moveJ} commands; the UR controller receives a
target joint configuration and drives the arm there using internally computed torques.
The payload mass parameter affects only these internal gravity-compensation and inertial
feedforward torques, not the commanded joint-angle profile.
No kinematic RMSE change from payload injection is architecturally possible under this control
scheme, and none is observed.
The result is a clean null finding on kinematic accuracy, consistent with the theoretical
prediction.


\section{TCP Path Deviation}
\label{sec:tcp_deviation}

\begin{figure}[H]
    \centering
    \includegraphics[width=0.88\textwidth]{Projekt/Figures/7_Test and Evaluation/fig8_tcp_deviation.pdf}
    \caption{TCP path deviation between the real robot and Pipeline~2 (\textit{sim\_aas})
             over the trajectory window (trajectory-onset aligned).}
    \label{fig:tcp_deviation}
\end{figure}

The TCP deviation results in Table~\ref{tab:tcp} are consistent with the joint RMSE results.
Angular errors of 30--37$^\circ$ at J3 and J4 propagate through the kinematic chain and
produce decimetre-scale displacement at the tool tip, visible in
Figure~\ref{fig:tcp_deviation}.

\begin{table}[htbp]
\centering
\caption{TCP RMS path deviation and its standard deviation over the trajectory window
         (trajectory-onset aligned).}
\label{tab:tcp}
\begin{tabular}{lcc}
\toprule
Comparison                          & TCP RMS deviation (mm) & Std (mm) \\
\midrule
real (P1) vs \textit{sim\_aas} (P2) & 225.9                  & 139.0 \\
\bottomrule
\end{tabular}
\end{table}

The RMS TCP deviation of 225.9\,mm reflects the same URSim motion-interpolation gap identified
in Section~\ref{sec:joint_accuracy}.
It is not related to the payload parameter and is not addressable through AAS injection
into URSim.


\section{Motor Current}
\label{sec:current_analysis}

Motor current reflects the dynamic model rather than just the kinematic trajectory.
A simulator that applies correct gravity compensation and inertial feedforward will draw current
patterns that match the physical robot, even if the joint positions are identical between
pipelines.
This makes current the most relevant metric for evaluating whether the AAS payload parameter
actually had an effect.

\begin{figure}[H]
    \centering
    \includegraphics[width=0.88\textwidth]{Projekt/Figures/7_Test and Evaluation/fig8_rms_current_bar.pdf}
    \caption{RMS motor current per joint for Pipeline~1 (\textit{real}) and Pipeline~2
             (\textit{sim\_aas}).
             Pipeline~2 closely matches the real robot at J2 and J3, the two joints that
             support the arm against gravity.
             Near-zero current at J1, J5, and J6 in the simulation reflects the absence of
             a friction model in URSim.}
    \label{fig:rms_current}
\end{figure}

\begin{table}[htbp]
\centering
\begin{tabular}{lrrr}
\toprule
Joint           & \textit{real} (P1) & \textit{sim\_aas} (P2) & Fidelity \\
\midrule
J1 -- Base      & 0.2767 & 0.0036 &   1\,\% \\
J2 -- Shoulder  & 1.2681 & 1.2558 &  99\,\% \\
J3 -- Elbow     & 1.1662 & 1.2028 & 103\,\% \\
J4 -- Wrist~1   & 0.6691 & 0.5688 &  85\,\% \\
J5 -- Wrist~2   & 0.1984 & 0.0109 &   5\,\% \\
J6 -- Wrist~3   & 0.2252 & 0.0008 & $<$1\,\% \\
\midrule
Mean            & 0.634  & 0.507  &  80\,\% \\
\bottomrule
\end{tabular}
\caption{RMS motor current (A) per joint. The fidelity column expresses Pipeline~2 current
         as a percentage of the real robot's current.}
\label{tab:current}
\end{table}

Two patterns stand out in Figure~\ref{fig:rms_current} and Table~\ref{tab:current}.

J1, J5, and J6 show near-zero current in Pipeline~2.
The physical robot draws measurable current at these joints even during slow phases, primarily
due to static holding torque and friction.
URSim does not model friction or back-EMF, so no AAS parameter can close this gap without a
corresponding simulator capability.
This is a known limitation of URSim as a program-validation tool rather than a
physics-accurate engine; it is precisely the gap that Pipeline~3 (the AAS Python Simulator,
which models friction via the \textit{DynamicsParameters} submodel) is designed to address.

AAS payload injection achieved high current fidelity at J2 and J3, the shoulder and elbow
joints that must counteract gravity across the entire suspended mass of the arm and tool.
Their current draw is therefore directly proportional to the total payload.
When Pipeline~2 registers the correct mass of 1.2\,kg, the simulator allocates proportionally
larger torques to these joints, and the resulting current closely matches the physical robot:

\begin{itemize}
    \item J2 (Shoulder): 1.256\,A versus 1.268\,A real, 99\,\% fidelity.
    \item J3 (Elbow): 1.203\,A versus 1.166\,A real, 103\,\% fidelity.
\end{itemize}

The overall mean current is 0.507\,A against the real robot's 0.634\,A (80\,\% overall
fidelity).
The residual deficit is concentrated entirely at J1, J5, and J6, where friction modelling
would be required regardless of AAS data quality.


\section{Summary of Findings}
\label{sec:eval_summary}

The experiment produces a \textbf{split outcome}: a null kinematic result and a clear dynamic
improvement.

On kinematic accuracy, the 21.3$^\circ$ combined RMSE and 225.9\,mm TCP RMS deviation reflect
the inherent gap between URSim's software-loop motion execution and the physical controller's
real-time interpolation.
This gap is unrelated to AAS parameter injection and cannot be reduced through payload injection
under \texttt{moveJ} joint-space control.
The null kinematic result is, in fact, strong confirmatory evidence that the alignment
methodology is sound and that the simulated trajectories are phase-coherent with the
physical robot's motion.

On dynamic accuracy, Pipeline~2 shows a clear and measurable improvement.
Injecting the correct payload mass brought J2 and J3 current to 99\,\% and 103\,\% of real
respectively.
This is the effect the AAS was designed to produce: providing an accurate, machine-readable
physical parameter to a connected system so that it can model the asset more faithfully.

Table~\ref{tab:summary} consolidates all metrics for the comparison between Pipeline~1 and
Pipeline~2.

\begin{table}[htbp]
\centering
\caption{Key evaluation metrics: Pipeline~1 (\textit{real}) versus Pipeline~2
         (\textit{sim\_aas}). Kinematic metrics use trajectory-onset alignment; current
         metrics are run-wide.}
\label{tab:summary}
\begin{tabular}{lccc}
\toprule
Metric                         & \textit{real} (P1) & \textit{sim\_aas} (P2) & Fidelity \\
\midrule
Mean cycle time (s)            & $34.67 \pm 0.02$   & $39.84 \pm 0.04$       & -- \\
Pre-trajectory dwell (s)       & $\approx$1.85      & $\approx$14.4          & -- \\
\midrule
Combined joint RMSE ($^\circ$) & --                 & 21.3                   & N/A$^{*}$ \\
TCP RMS deviation (mm)         & --                 & 225.9                  & N/A$^{*}$ \\
\midrule
Mean RMS current (A)           & 0.634              & 0.507                  & 80\,\% \\
J2 current (A)                 & 1.268              & 1.256                  & 99\,\% \\
J3 current (A)                 & 1.166              & 1.203                  & 103\,\% \\
\bottomrule
\multicolumn{4}{l}{\footnotesize $^{*}$Kinematic accuracy is constrained by URSim's motion
interpolation, not by AAS parameters;} \\
\multicolumn{4}{l}{\footnotesize \phantom{$^{*}$}no improvement is architecturally
possible under \texttt{moveJ} control (see Section~\ref{sec:joint_accuracy}).} \\
\end{tabular}
\end{table}

The residual current gap at J1, J5, and J6 reflects the absence of friction modelling in
URSim, a limitation that falls outside the parameter surface available to Pipeline~2.
Addressing this gap is the stated purpose of Pipeline~3, which is designed to inject the full
\textit{DynamicsParameters} submodel into a custom Python simulator capable of modelling
friction effects directly.
The implications of this split outcome for the research question are discussed in
Chapter~\ref{chap:discussion}.
